\documentclass[a4paper,10pt]{article} \usepackage{anysize}
\marginsize{2cm}{2cm}{1cm}{1cm}
%\textwidth 6.0in \textheight = 664pt
\usepackage{xltxtra}
\usepackage{graphicx}
\usepackage{color}
\usepackage{xgreek}
\usepackage{fancyvrb}
\usepackage{minted}
\usepackage{listings}
\usepackage{enumitem}
\usepackage{framed}
\usepackage{relsize}
\usepackage{float}
\usepackage{amsmath}
\usepackage{mathtools}
\setmainfont[Mapping=TeX-text]{CMU Serif}

%\setmainfont{Arial}
%\setsansfont{FreeSans}
%\setmonofont{FreeMono}
\begin{document}
%\def\thesubsection {\alph{subsection}}
\renewcommand{\labelenumi}{\roman{enumi})}
\renewcommand{\labelenumii}{(\arabic{enumii})}

\begin{titlepage}
\begin{center}
\begin{figure}[t] 
     \includegraphics[scale=0.7]{title/ntua_logo}
\end{figure}
\begin{LARGE}\textbf{ΕΘΝΙΚΟ ΜΕΤΣΟΒΙΟ ΠΟΛΥΤΕΧΝΕΙΟ\\}\end{LARGE}
\vspace{5cm}
\begin{Large}
ΣΧΟΛΗ ΗΜ\&ΜΥ\\
Αλγόριθμοι και Πολυπλοκότητα\\
1\textsuperscript{η} Σειρά Γραπτών Ασκήσεων\\
Ακ. έτος 2010-2011\\
\end{Large}
\vspace{10cm}
\begin{tabular}{l r}
\Large{Λύρας Γρηγόρης}&
\large{Α.Μ.: 03109687}\\
\end{tabular}\\
\vspace{1cm}

\vfill
\large\today\\
\end{center}
\end{titlepage}



\section{Επιτροπή Αντιπροσώπων} \setcounter{section}{1}
Κάνουμε ταξινόμηση στον πίνακα των $[s_i,f_i)$ ως προς $f_i$. Ξεκινώντας από
 το πρώτο $(n=0)$, ελέγχω ποια από τα επόμενα στον ταξινομημένο πίνακα επικαλύπτονται
με αυτό μέχρι να βρω κάποιο με το οποίο δεν έχει επικάλυψη έστω $[s_k,f_k)$. Κρατάω
το τελευταίο $([s_{k-1},f_{k-1}))$ από τα προηγούμενα. Στη συνέχεια, αν αυτό έχει
επικάλυψη με το $[s_k,f_k)$ συνεχίζω για $n=k+1$. Αλλιώς συνεχίζω
τον ίδιο αλγόριθμο για $n=k$.


\vspace{3cm}

\section{Βιαστικός Μοτοσικλετιστής}
Ταξινομούμε τον πίνακα των ταχυτήτων-αποστάσεων με βάση τις ταχύτητες.
Ξεκινώντας από τη μικρότερη "τρέχουμε" στο αντίστοιχο διάστημα με τη μέγιστη
ταχύτητα ξεπερνώντας και το όριο κατά $u$. Αυτό μας παίρνει χρόνο
$T_i=\frac{l_i}{u_i+u}$. Αφαιρούμε το χρόνο αυτό από τον συνολικό χρόνο που
έχουμε και συνεχίζουμε (προφανώς αν ξεπεράσουμε το χρονικό όριο η απόσταση που
θα διανύσουμε μέσα σε αυτό το τμήμα θα είναι $x=l_i*\frac{T}{T_i}$), για την
επόμενη σε σειρά ταχύτητα. Με τον ίδιο τρόπο. Μόλις ο χρόνος που έχουμε
εξαντληθεί, σταματάμε. Τα τμήματα που έχουμε επιλέξει για να υπερβούμε το όριο
ταχύτητας θα ελαχιστοποιήσουν το χρόνο άφιξης στο $B$.

Έστω $T_s$ ο συνολικός χρόνος από το $A$ στο $B$. Συνεπώς έχουμε τη σχέση:
\[
T_s=\sum_{i=1}^n{\frac{l_i}{u_i}}
\]
Θέλουμε να ελαχιστοποιήσουμε το άθροισμα αυτό από θετικούς όρους (η ταχύτητα
δεν έχει νόημα να είναι προσημασμένη σε αυτό το πρόβλημα), συνεπώς θέλουμε να
ελαχιστοποιήσουμε τους όρους $k_1, k_2, k_3 ...$.

Έστω ο όρος $k$:
\[
T_k=
	\begin{dcases}
	\frac{l_i}{u_i} \\
	\frac{l_i}{u_i+u}  
	\end{dcases}
\]

Ορίζω $\Delta{T}$ τη διαφορά:
\[
	\Delta{T}=\frac{l_i}{u_i} - \frac{l_i}{u_i+u} = l_i*\frac{u}{u_i*(u_i+u)}
\] 
\[
	\frac{\theta{\Delta{T}}}{\theta{u_i}} =
	l_i*u*(-\frac{2*u_i+u}{(u_i*(u_i+u))^2}) < 0
\]

Συνεπώς αυξάνοντας το $u_i$ μειώνεται το $\Delta{T}$, το κέρδος δηλαδή που θα
είχαμε ξεπερνώντας το όριο ταχύτητας. Αυτό συνεπάγεται πως πρώτα θέλουμε να
πάρουμε τα κομμάτια με τη μικρότερη ταχύτητα καθώς αυτά μας δίνουν το μέγιστο
κέρδος. Το πρόβλημα λοιπόν ανάγεται σε συνεχές Knapsack όπου το Greedy
κριτήριο είναι η "μικρότερη ταχύτητα πρώτη".

Σε περίπτωση που αποφασίσουμε να υπερβούμε την ταχύτητα κατά έναν παραγοντα
$a>1$ τότε η $\Delta{T}$ έχει τη μορφή:
\[
	\Delta{T}=\frac{l_i}{u_i} - \frac{l_i}{a*u_i} =
	\frac{a-1}{a}*\frac{l_i}{u_i} = \frac{a-1}{a}*T_i
\] 

Συνεπώς τώρα το κριτήριο αλλάζει και πλέον θέλουμε πρώτα τα κομμάτια που
παίρνουν περισσότερο χρόνο για να τα διανύσουμε. Κατά τα άλλα ο αλγόριθμος
παραμένει ίδιος.





\vspace{3cm}

\section{Βότσαλα στη Σκακιέρα}
Ορίζω τη σχέση:
\[
	dTable[i]:[k_1],[k_2],..,[k] = dTable[i][k_1],dTable[i][k_2],..,dTable[i][k]
\]
\begin{tabular}{| c | l | l  c  l |}
\hline
&iteration $i$&iteration $i+1$&&\\
\hline
$X+Z$ &dTable[i][0]&$X[i+1]+Z[i+1]$&$+$&$max\{dTable[i]:[1],[4],[6],[7]\}$ 
\\
$Y+W$ &dTable[i][1]&$Y[i+1]+W[i+1]$&$+$&$max\{dTable[i]:[0],[3],[5],[7]\}$  
\\
$X+W$ &dTable[i][2]&$X[i+1]+W[i+1]$&$+$&$max\{dTable[i]:[4],[5],[7]\} $
\\
$X$ &dTable[i][3]&$X[i+1]$&$+$&$max\{dTable[i]:[1],[4],[5],[6],[7]\} $
\\
$Y$ &dTable[i][4]&$Y[i+1]$&$+$&$max\{dTable[i]:[0],[2],[3],[5],[6],[7]\} $
\\
$Z$ &dTable[i][5]&$Z[i+1]$&$+$&$max\{dTable[i]:[1],[2],[3],[4],[6],[7]\} $
\\
$W$ &dTable[i][6]&$W[i+1]$&$+$&$max\{dTable[i]:[0],[3],[4],[5],[7]\} $
\\
$0$ &dTable[i][7]&$0$&$+$&$max\{dTable[i]:[0],[1],[2],[3],[4],[5],[6],[7]\} $
\\

\hline
\end{tabular}



\vspace{3cm}

\section{Χωρισμός Κειμένου σε Γραμμές}


\vspace{3cm}

\section{Αντίγραφα Αρχείου}


\vspace{3cm}

\section{Έλεγχος Ταξινόμησης}


\vspace{3cm}

\section{Bonus: Δρομολόγηση Εργασιών}


\end{document}
